% year: תשפ"ה
% type: מבחן
% semester: א
% moed: א
% number: 1ב
% points: 10
% categories: ivt, func-analysis
% source: exams/מבחנים/תשפה/מועד א תשפה.pdf

\begin{question}
(10 נק') כמה פתרונות ממשיים יש למשוואה הבאה?
\[ x + \sqrt{x} = 17 \]
נמקו היטב.
\end{question}

\begin{hint}
הגדירו $g(x) = x + \sqrt{x} - 17$ על $[0,\infty)$: קיום פתרון נובע ממשפט ערך הביניים, ויחידותו ממונוטוניות.
\end{hint}

\begin{solution}
נגדיר $g(x) = x + \sqrt{x} - 17$. המשוואה מוגדרת רק עבור $x \ge 0$, ו-$g$ רציפה על $[0,\infty)$.

\textbf{קיום:} $g(0) = -17 < 0$ ו-$g(17) = \sqrt{17} > 0$. לפי משפט ערך הביניים קיימת $c \in (0,17)$ עם $g(c) = 0$.

\textbf{יחידות:} עבור $x > 0$ מתקיים $g'(x) = 1 + \frac{1}{2\sqrt{x}} > 0$, ולכן (מסקנה ממשפט לגרנז') $g$ עולה ממש על $[0,\infty)$
(למעשה זה ברור גם ישירות: סכום של שתי פונקציות עולות ממש). פונקציה עולה ממש מקבלת כל ערך לכל היותר פעם אחת, ולכן יש לכל היותר פתרון אחד.

\textbf{תשובה:} למשוואה פתרון ממשי \textbf{אחד בדיוק}. (למעשה, בהצבה $t = \sqrt{x}\ge 0$ מקבלים $t^2 + t - 17 = 0$, כלומר $t = \frac{-1+\sqrt{69}}{2}$, ו-$x = \left(\frac{-1+\sqrt{69}}{2}\right)^2 \approx 13.35$.)
\end{solution}
